\begin{theorem}
  \label{ent:thm:floor-spec}
  For every real $x$, its greatest integer satisfies $\lfloor x\rfloor\le x<\lfloor x\rfloor+1$.
  In particular, $x=\lfloor x\rfloor+\theta$ for a real $0\le\theta<1$.

  \begin{lean}
    theorem floor_spec (x : ℝ) :
        (⌊x⌋ : ℝ) ≤ x ∧ x < (⌊x⌋ : ℝ) + 1 ∧
          ∃ θ : ℝ, 0 ≤ θ ∧ θ < 1 ∧ x = ⌊x⌋ + θ
  \end{lean}
\end{theorem}

\begin{proof}
  The defining floor bounds hold, and $\theta=x-\lfloor x\rfloor$ gives the decomposition.
  Its bounds follow by subtraction; its value is forced by the displayed equality.

  \begin{lean}
    have hlo := Int.floor_le x

    have hhi := Int.lt_floor_add_one x

    exact ⟨hlo, hhi, x - ⌊x⌋, by linarith, by linarith, by ring⟩
  \end{lean}
\end{proof}

\begin{theorem}
  \label{ent:thm:floor-eq-self-iff}
  A real number equals its floor exactly when it is an integer.

  \begin{lean}
    theorem floor_eq_self_iff (x : ℝ) : (⌊x⌋ : ℝ) = x ↔ ∃ z : ℤ, x = z
  \end{lean}
\end{theorem}

\begin{proof}
  Equality with the floor supplies the integer witness.
  Conversely, the floor of an integer is that integer.

  \begin{lean}
    constructor

    · intro h
      exact ⟨⌊x⌋, h.symm⟩

    · rintro ⟨z, rfl⟩
      rw [Int.floor_intCast]
  \end{lean}
\end{proof}

\begin{theorem}
  \label{ent:thm:floor-nat-div}
  For natural $N$ and positive $k$, the integer quotient is $N/k=\lfloor N/k\rfloor$, where the
  quotient inside the floor is taken in the real numbers.

  \begin{lean}
    theorem floor_nat_div (N k : ℕ) (hk : 0 < k) :
        ⌊(N : ℝ) / k⌋ = ((N / k : ℕ) : ℤ)
  \end{lean}
\end{theorem}

\begin{proof}
  The division identity $N=kq+r$ with $0\le r<k$ gives $q\le N/k<q+1$.
  These bounds identify the floor with $q$.

  \begin{lean}
    have hkreal : (0 : ℝ) < k := by exact_mod_cast hk
    have heq := Nat.mod_add_div N k
    have hlt := Nat.mod_lt N hk

    apply Int.floor_eq_iff.mpr
    constructor

    · apply (le_div_iff₀ hkreal).mpr
      exact_mod_cast Nat.div_mul_le_self N k

    · apply (div_lt_iff₀ hkreal).mpr
      simp only [Int.cast_natCast]

      have hreal : (N % k : ℕ) + (k : ℝ) * (N / k : ℕ) = N := by exact_mod_cast heq
      have hrlt : ((N % k : ℕ) : ℝ) < k := by exact_mod_cast hlt
      linarith
  \end{lean}
\end{proof}

\begin{theorem}
  \label{ent:thm:card-multiples}
  The number of multiples of positive $k$ in $\{1,\ldots,N\}$ is $\lfloor N/k\rfloor$.

  \begin{lean}
    theorem card_multiples (N k : ℕ) (hk : 0 < k) :
        ((Finset.Icc 1 N).filter (fun n => k ∣ n)).card = N / k
  \end{lean}
\end{theorem}

\begin{proof}
  Send a multiple $n$ to $n/k$.
  This is a bijection with $\{1,\ldots,\lfloor N/k\rfloor\}$, whose inverse is multiplication by
  $k$.

  \begin{lean}
    classical

    have hcard : ((Finset.Icc 1 N).filter (fun n => k ∣ n)).card =
        (Finset.Icc 1 (N / k)).card := by

      apply Finset.card_bij (fun n _ => n / k)

      · intro n hn

        obtain ⟨hrange, hdiv⟩ := Finset.mem_filter.mp hn
        obtain ⟨hnpos, hnle⟩ := Finset.mem_Icc.mp hrange

        apply Finset.mem_Icc.mpr
        exact ⟨Nat.div_pos (Nat.le_of_dvd hnpos hdiv) hk, Nat.div_le_div_right hnle⟩

      · intro m hm n hn heq

        have hmd := (Finset.mem_filter.mp hm).2

        have hnd := (Finset.mem_filter.mp hn).2

        have h := congrArg (fun q => q * k) heq
        rwa [Nat.div_mul_cancel hmd, Nat.div_mul_cancel hnd] at h

      · intro j hj

        obtain ⟨hjpos, hjle⟩ := Finset.mem_Icc.mp hj

        refine ⟨j * k, ?_, Nat.mul_div_cancel j hk⟩

        apply Finset.mem_filter.mpr
        constructor

        · apply Finset.mem_Icc.mpr
          exact ⟨mul_pos hjpos hk, (Nat.le_div_iff_mul_le hk).1 hjle⟩

        · exact ⟨j, by ring⟩

    rw [hcard, Nat.card_Icc, Nat.add_sub_cancel]
  \end{lean}
\end{proof}

\begin{theorem}
  \label{ent:thm:legendre}
  For a prime $p$ and natural $n$, the exponent of $p$ in $n!
  $ is \[ v_p(n!)=\sum_{j=1}^{n}\left\lfloor\frac{n}{p^j}\right\rfloor. \] Larger indices contribute zero, so this is also the usual finite sum over $j\ge1$.

  \begin{lean}
    theorem legendre (n p : ℕ) (hp : Chapter02.isPrime p) :
        primeExponents n.factorial p = ∑ j ∈ Finset.Icc 1 n, n / p ^ j
  \end{lean}
\end{theorem}

\begin{proof}
  In the checked bounded formula, sum $v_p(k)$ over $1\le k\le n$.
  The exponent $v_p(k)$ counts the positive powers of $p$ dividing $k$.
  Interchange these finite sums.
  For each $j$, the inner count is the number of multiples of $p^j$ up to $n$, namely $\lfloor
    n/p^j\rfloor$.
  The bound $v_p(k)\le k$ justifies truncation.

  \begin{lean}
    legendre_with_bound n p n hp le_rfl
  \end{lean}
\end{proof}

\begin{theorem}
  \label{ent:thm:floor-add-carry}
  For real $a,b$, the difference $\lfloor a+b\rfloor-\lfloor a\rfloor-\lfloor b\rfloor$ is either
  zero or one.

  \begin{lean}
    theorem floor_add_carry (a b : ℝ) :
        ⌊a + b⌋ - ⌊a⌋ - ⌊b⌋ = 0 ∨ ⌊a + b⌋ - ⌊a⌋ - ⌊b⌋ = 1
  \end{lean}
\end{theorem}

\begin{proof}
  Add the floor bounds for $a$ and $b$ and compare with those for $a+b$.
  The integer difference is strictly greater than minus one and strictly less than two, leaving
  exactly the stated possibilities.

  \begin{lean}
    have ha := floor_spec a

    have hb := floor_spec b

    have hab := floor_spec (a + b)

    have hlow : (⌊a⌋ : ℝ) + ⌊b⌋ - 1 < ⌊a + b⌋ := by linarith

    have hhigh : (⌊a + b⌋ : ℝ) < ⌊a⌋ + ⌊b⌋ + 2 := by linarith

    have hlow' : ⌊a⌋ + ⌊b⌋ - 1 < ⌊a + b⌋ := by exact_mod_cast hlow

    have hhigh' : ⌊a + b⌋ < ⌊a⌋ + ⌊b⌋ + 2 := by exact_mod_cast hhigh

    omega
  \end{lean}
\end{proof}

\begin{theorem}
  \label{ent:thm:floor-add-ge}
  For all real $a,b$, $\lfloor a+b\rfloor\ge\lfloor a\rfloor+\lfloor b\rfloor$.

  \begin{lean}
    theorem floor_add_ge (a b : ℝ) : ⌊a⌋ + ⌊b⌋ ≤ ⌊a + b⌋
  \end{lean}
\end{theorem}

\begin{proof}
  The difference is zero or one, hence is nonnegative.

  \begin{lean}
    rcases floor_add_carry a b with h | h <;> omega
  \end{lean}
\end{proof}

\begin{theorem}
  \label{ent:thm:factorials-dvd-factorial}
  For natural $r\le n$, the factorial quotient $n!
    /(r!(n-r)!)$ is an integer.
  This includes the endpoint cases $r=0,n$, when the quotient is one.

  \begin{lean}
    theorem factorials_dvd_factorial (n r : ℕ) (hr : r ≤ n) :
        r.factorial * (n - r).factorial ∣ n.factorial
  \end{lean}
\end{theorem}

\begin{proof}
  It suffices to compare every prime exponent in the denominator and numerator.
  Legendre\textquotesingle s formula with a common upper bound reduces this to \[ \left\lfloor\frac r{p^j}\right\rfloor+\left\lfloor\frac{n-r}{p^j}\right\rfloor\le\left\lfloor\frac n{p^j}\right\rfloor. \] Multiplying the two integer quotients by $p^j$ shows that their sum is at most $n/p^j$.
  Sum these inequalities and apply the exponent criterion for divisibility.

  \begin{lean}
    apply (dvd_iff_primeExponents_le _ _
      (mul_pos (Nat.factorial_pos r) (Nat.factorial_pos (n - r))) (Nat.factorial_pos n)).2

    intro p
    by_cases he : primeExponents (r.factorial * (n - r).factorial) p = 0

    · rw [he]
      exact Nat.zero_le _

    · have hp := (primeExponents_spec (r.factorial * (n - r).factorial)
        (mul_pos (Nat.factorial_pos r) (Nat.factorial_pos (n - r)))).1 p
        (Finsupp.mem_support_iff.mpr he)

      rw [primeExponents_mul _ _ (Nat.factorial_pos r) (Nat.factorial_pos (n - r)),
        Finsupp.add_apply, legendre_with_bound r p n hp hr,
        legendre_with_bound (n - r) p n hp (Nat.sub_le n r), legendre n p hp,
        ← Finset.sum_add_distrib]

      apply Finset.sum_le_sum
      intro j hj

      have hpos : 0 < p ^ j := pow_pos (by
          have := hp.1
          omega) j

      apply (Nat.le_div_iff_mul_le hpos).2

      have h₁ := Nat.div_mul_le_self r (p ^ j)

      have h₂ := Nat.div_mul_le_self (n - r) (p ^ j)

      have hsum : r + (n - r) = n := Nat.add_sub_cancel' hr

      nlinarith
  \end{lean}
\end{proof}

\begin{theorem}
  \label{ent:thm:factorial-dvd-consecutive-product}
  The product of any $r$ consecutive positive integers is divisible by $r!
  $.
  More precisely, $r!
    \mid(a+1)(a+2)\cdots(a+r)$ for natural $a,r$.

  \begin{lean}
    theorem factorial_dvd_consecutive_product (a r : ℕ) :
        r.factorial ∣ ∏ k ∈ Finset.Ico (a + 1) (a + r + 1), k
  \end{lean}
\end{theorem}

\begin{proof}
  Multiplication of the displayed product by $a!
  $ gives $(a+r)!$.
  The preceding theorem gives $r!
    a!\mid(a+r)!$.
  Cancel the positive factor $a!
  $ to obtain the claimed divisibility.

  \begin{lean}
    have hfactor : (∏ k ∈ Finset.Ico (a + 1) (a + r + 1), k) * a.factorial =
        (a + r).factorial := by
      induction r with
      | zero => simp only [Nat.add_zero, Finset.Ico_self, Finset.prod_empty, one_mul]

      | succ r ih =>
        rw [show a + (r + 1) + 1 = (a + r + 1) + 1 by omega,
          Finset.prod_Ico_succ_top (by omega : a + 1 ≤ a + r + 1)]
        rw [mul_right_comm, ih]
        rw [show a + (r + 1) = (a + r) + 1 by omega, Nat.factorial_succ]
        ring

    obtain ⟨k, hk⟩ := factorials_dvd_factorial (a + r) r (by omega)
    rw [show a + r - r = a by omega] at hk

    refine ⟨k, ?_⟩
    apply Nat.eq_of_mul_eq_mul_right (Nat.factorial_pos a)
    rw [hfactor, hk]
    ring
  \end{lean}
\end{proof}

\begin{theorem}
  \label{ent:thm:sum-divisorsum}
  For $F(n)=\sum_{d\mid n}f(d)$ and every natural $N$, \[ \sum_{n=1}^N F(n)=\sum_{k=1}^N f(k)\left\lfloor\frac Nk\right\rfloor. \] No multiplicativity assumption is needed.

  \begin{lean}
    theorem sum_divisorSum (f : ℕ → ℤ) (N : ℕ) :
        (∑ n ∈ Finset.Icc 1 N, divisorSum f n) =
          ∑ k ∈ Finset.Icc 1 N, f k * (N / k : ℕ)
  \end{lean}
\end{theorem}

\begin{proof}
  Embed every divisor sum in the common index interval $1,\ldots,N$, assigning zero to
  nondivisors.
  Interchange the finite sums.

  \begin{lean}
    classical

    have hexpand (n : ℕ) (hn : n ∈ Finset.Icc 1 N) :
        divisorSum f n = ∑ d ∈ Finset.Icc 1 N, if d ∣ n then f d else 0 := by
      obtain ⟨hnpos, hnle⟩ := Finset.mem_Icc.mp hn

      have hset : n.divisors = (Finset.Icc 1 N).filter (fun d => d ∣ n) := by
        ext d
        constructor

        · intro hd
          exact Finset.mem_filter.mpr
            ⟨Finset.mem_Icc.mpr ⟨Nat.pos_of_mem_divisors hd,
              (Nat.le_of_dvd hnpos (Nat.dvd_of_mem_divisors hd)).trans hnle⟩,
              Nat.dvd_of_mem_divisors hd⟩

        · intro hd
          exact Nat.mem_divisors.mpr ⟨(Finset.mem_filter.mp hd).2, ne_of_gt hnpos⟩

      rw [divisorSum, hset, Finset.sum_filter]
  \end{lean}

  For a fixed $k$, the term $f(k)$ occurs once for each multiple of $k$ up to $N$.
  The multiple-counting theorem makes this number $\lfloor N/k\rfloor$.

  \begin{lean}
    rw [Finset.sum_congr rfl hexpand, Finset.sum_comm]
    apply Finset.sum_congr rfl
    intro k hk

    rw [← Finset.sum_filter, Finset.sum_const, nsmul_eq_mul,
      card_multiples N k (Finset.mem_Icc.mp hk).1, mul_comm]
  \end{lean}
\end{proof}

\begin{theorem}
  \label{ent:thm:sum-tau}
  For natural $N$, \[ \sum_{n=1}^N\tau(n)=\sum_{k=1}^N\left\lfloor\frac Nk\right\rfloor. \]

  \begin{lean}
    theorem sum_tau (N : ℕ) :
        (∑ n ∈ Finset.Icc 1 N, (tau n : ℤ)) = ∑ k ∈ Finset.Icc 1 N, (N / k : ℕ)
  \end{lean}
\end{theorem}

\begin{proof}
  Apply the summatory divisor-sum identity to the constant function one.

  \begin{lean}
    have h := sum_divisorSum (fun _ => 1) N

    simpa only [divisorSum, Finset.sum_const, nsmul_eq_mul, mul_one, one_mul, tau, Nat.cast_sum] using h
  \end{lean}
\end{proof}

\begin{theorem}
  \label{ent:thm:sum-sigma}
  For natural $N$, \[ \sum_{n=1}^N\sigma(n)=\sum_{k=1}^N k\left\lfloor\frac Nk\right\rfloor. \]

  \begin{lean}
    theorem sum_sigma (N : ℕ) :
        (∑ n ∈ Finset.Icc 1 N, (sigma n : ℤ)) =
          ∑ k ∈ Finset.Icc 1 N, (k : ℤ) * (N / k : ℕ)
  \end{lean}
\end{theorem}

\begin{proof}
  Apply the same identity to $f(k)=k$.

  \begin{lean}
    have h := sum_divisorSum (fun k => (k : ℤ)) N

    simpa only [divisorSum, sigma, Nat.cast_sum] using h
  \end{lean}
\end{proof}
